\section{Calculation by Aagaard}
\begin{multicols}{2}
    \subsection*{1. Romanko - Shulakova}
    \newchessgame[
        setfen=r4rk1/pb3pp1/1p1ppq1p/8/2PR3n/P1N1P1P1/1PQ1BP1P/5RK1 b,
        moveid=1b
    ]
    \chessboard

    \mythoughts{
        Clearly Black has something along the long diagonal. 
        \variation{
            1... Nf3+ 2. Kg2 Nxd4+ 3. Kg1 Nxc2
        } Black wins a lot of material.
        White can however improve his play:
        \variation{
            1... Nf3+ 2. Bxf3 Qxf3 3. e4
        } Black gains nothing. 

        Often it is a good idea to switch the order of the moves. Let's try:
        \mainline{
            1... Qf3 2. Bxf3 Nxf3+ 3. Kg2 Nxd4+ 4. Kg1 Nxc2
        } Black wins material so that's the move!
    }

    \subsection*{2. Hunt - Galdunts}
    \newchessgame[
        setfen=1br3k1/5pp1/p3p2p/1bBn4/N7/1P3PP1/5R1P/2R3K1 b,
        moveid=1b
    ]
    \chessboard

    \mythoughts{
        White's pieces are not well placed. His king, rook and bishop are on the same diagonal.
        His bishop and another rook are pinned on the same file. So there should be 
        some pin here.  

        The knight on a4 should be removed as soon as possible since it protects the bishop on c5.

        \mainline{
            1... Bxa4 2. bxa4 Ba7 3. Rfc2 
        } are natural moves.
    }

    \chessboard

    \mythoughts{
        Going further: \mainline{3... Nb4 } Black threatens to take the rook on c5
        and fork with his knight on d3. Black wins material.
    }

    \subsection*{3 Baklan - Getz}
    \newchessgame[
        setfen=4rrk1/R5pp/3qp3/1p1b1p1P/5n2/2P1QN2/1PB2PP1/3R2K1 w,
        moveid=1w
    ]
    \chessboard
\end{multicols}